\documentclass[main.tex]{subfiles}
\begin{document}
%% All commands are in macros.sty

\chapter{Supersingular Reduction of Hurwitz Curves}

\section{Introduction}

Let $C / K$ be a smooth projective curve over a field of characteristic zero. Hurwitz proved a uniform bound on the number of automorphisms of $C$ depending on the genus $g$,
\[ \# \Aut(C) \le 84(g-1). \]
We call a curve achieving this bound a \textit{Hurwitz curve}. Every Hurwitz curve is defined over a number field because the quotient by the automorphism group $C \to C / \Aut(C) \cong \P^1$ canonically presents $C$ as a $G$-\Belyi map with orders of monodromy $(2,3,7)$ over $(0,1,\infty)$. The arithmetic of Hurwitz curves is quite rich, for example Gonz\'alez-Diez and Jaikin-Zapirain \cite{gonzalez-diez:galois_moduli} showed that the action of $\Gal(\ol{\Q}/\Q)$ on the set of isomorphism classes of Hurwitz curves is faithful. However, Hurwitz curves only appear in a sparse set of genera. Their geometric classification can be entirely reduced to group theory. Indeed, over $\CC$, the set of isomorphism classes of pairs $(C, \varphi : \Aut(C) \iso G)$ where $C$ is a Hurwitz curve and $\varphi : \Aut(C) \iso G$ is an identification of its automorphism group with $G$ is in canonical bijection with the set of surjections
\[ \Delta(2,3,7) \onto G \]
up to conjugacy where $\Delta(a,b,c)$ is the von Dyck group
\[ \Delta(a,b,c) := \left< x,y,z \mid x^a, y^b, z^c, xyz \right> \]
which is the index $2$ subgroup consisting of orientation-preserving isometries of the $(a,b,c)$-hyperbolic triangle group. Then $\Out(G)$ acts freely on the set of conjugacy classes of surjections whose orbits correspond to unmarked Hurwitz curves.

\begin{figure}
\begin{center}
    \begin{tabular}{c|c|c|c}
        $g$ & $G$ & $N$ & $\Out(G)$ \\
        \hline
         $3$ & $\PSL_2(\FF_7)$ & $1$ & $\Z / 2 \Z$
         \\
         $7$ & $\SL_2(\FF_8)$ & $1$ & $\Z / 3 \Z$
         \\
         $14$ & $\PSL_2(\FF_{13})$ & $3$ & $\Z / 2 \Z$
         \\
         $17$ & $G_{17}$ & $2$ & $\Z / 2 \Z$
    \end{tabular}
\end{center}
    \caption{The Hurwitz curves of genus $g \le 100$. Here $G$ is the automorphism group (of which there is a unique isomorphism class for each genus in the table), $N$ is the number of non-isomorphic geometric curves of genus $g$, and $\Out(G)$ is the outer automorphism group of $G$ \cite{Conder:survey, Conder:update}.}
    \label{table:low_genus}
\end{figure}

Hurwitz curves of genus $\le 100$ have been tabulated by Conder \cite{Conder:survey, Conder:update} and their results are summarized in table~\ref{table:low_genus}. Here $G_{17}$ is the non-split extension $2^3.\PSL_2(\FF_7)$ of order $1344$, arising from Macbeath's homology cover construction \cite{Macbeath:eightfold, macbeath:hurwitz_theorem} applied to the Klein quartic. It fits in a non-split exact sequence $0 \to (\Z/2\Z)^3 \to G_{17} \to \PSL_2(\FF_7) \to 0$ corresponding to an irreducible $3$-dimensional $\FF_2$-representation of $\PSL_2(\FF_7)$. The fact that $\mathrm{(P)SL}_2$ of a finite field shows up as the automorphism group for $g \le 14$ is not a coincidence. Indeed, those curves are all \textit{Shimura curves} for a particular quaternion algebra over the totally real field $F = \Q(\zeta_7)^+$ (see \S\ref{sec:shimura}). However, the curves of genus $17$ are uniformized by a \textit{non-congruence} normal subgroup of $\Delta(2,3,7)$ and hence do not arise as Shimura curves.

We say an abelian variety $A / k$ over a field of characteristic $p$ is \textit{supersingular} if it is isogenous to $E^g$ where $E$ is a supersingular elliptic curve over $\bar{k}$ \cite{Oort:subvarieties}. Over a finite field $k = \FF_q$, this is equivalent to requiring that the eigenvalues of $\Frob_q$ on $H_{\et}^1(A_{\bar{k}}, \Q_{\ell})$ are all of the form $\zeta \cdot q^{1/2}$ for $\zeta$ a root of unity \cite{tate:endomorphisms}. Over a general field of characteristic $p$, the definition is geometric: $A$ is supersingular if $A_{\bar{k}}$ is isogenous to a power of a supersingular elliptic curve.

We say a smooth curve $C / k$ is \textit{supersingular} if $\Jac(C)$ is supersingular. It is unknown what pairs $(p,g)$ appear over primes $p$ so that there is a supersingular curve of genus $g$ in characteristic $p$.

Given a curve $C / K$ over a number field, it is very interesting to ask about the set of primes of $K$ for which the reduction of $C$ is supersingular. This set is often sparse and, in large genus, expected to be finite. Elkies \cite{elkies:supersingular} proved that if $E / K$ is an elliptic curve and $K$ has a real place then $E$ has infinitely many primes of supersingular reduction. We prove that Hurwitz curves of small genera have infinitely many primes of supersingular reduction.

\begin{thm}\label{thm:main}
Let $C/K$ be a Hurwitz curve of genus $g = 3,7,14,17$. Then there are infinitely many primes of $K$ at which $C$ has supersingular reduction.
\end{thm}

This will follow from a careful analysis of the representation theory of their automorphism group and the arithmetic of certain quotients. The proof is given genus by genus and is spread over \S\ref{sec:genus3_7}--\S\ref{sec:genus17}. One motivation for this result is the author's recent counterexample to Shioda's conjecture on unirationality for surfaces. This conjecture, due to Shioda \cite{Shioda:unirational}, predicts that a smooth projective surface $X$ with $\pi_1^{\et}(X) = 1$ over an algebraically closed field of characteristic $p > 0$ is unirational if and only if it is supersingular.

To construct further examples, one needs a source of high genus curves with large automorphism groups admitting many primes of supersingular reduction. The Hurwitz curves therefore provide natural candidates. Theorem~\ref{thm:main} in $g = 14,17$ can be used to produce two infinite lists of primes at which Shioda's conjecture is known to fail. These lists are ultimately related to the supersingular primes of some particular elliptic curves. Hence the number of these primes up to $x$ are expected to satisfy the Lang--Trotter estimates $\sim x^{1/2 - o(1)}$ \cite{lang:frobenius_distributions}.

\section{Shimura curves} \label{sec:shimura}

We survey Shimura's groundbreaking work \cite{shimura:curves} constructing canonical models of Shimura curves with an eye toward the arithmetic of certain Hurwitz curves. The main purpose of this section is to prove that the elliptic curve $E$ appearing as an isogeny factor in the Jacobian of a genus $14$ Hurwitz curve $C$ is defined over a number field with a real place; see Theorem~\ref{app:thm:isogeny_decomp_defined_over_K}. This is necessary to apply Elkies' theorem \cite{elkies:supersingular} to show that $C$ has infinitely many primes of supersingular reduction.

\subsection{Canonical models of Shimura curves}

The exposition follows \S2--3 of \cite{shimura:curves}. Let $F$ be a totally real number field and $B$ a quaternion algebra over $F$. Write $g = [F : \Q]$ and $\frD(B/F)$ for the product of the non-archimedean primes of $F$ at which $B$ is \textit{non-split}, meaning $B \ot_F F_{\p}$ is a division algebra over $F_{\p}$. % REVIEW: "split" and "ramified" are opposite notions for quaternion algebras; non-split places are the ramified ones. Check which term is intended here.
At all other finite places $B$ is \textit{split} (or \textit{ramified}) meaning $B \ot_F F_{\p} \cong M_2(F_{\p})$. Let $\nu_1, \dots, \nu_g : F \embed \RR$ be the archimedean places ordered such that $B$ is split at $\nu_1, \dots, \nu_r$ and non-split at $\nu_{r+1}, \dots, \nu_{g}$. We also choose an isomorphism compatible with these fixed places
\[ B \ot_{\QQ} \RR \iso M_2(\RR)^r \times \HH^{g-r} \]
where $\HH$ are the Hamiltonian quaternions. Assuming $r > 0$, this defines an embedding
\[ \iota_{\infty} : B \embed M_2(\RR)^r. \]
Let $B^+$ denote the elements of $B$ with totally positive reduced norm (note that Shimura writes $N_{B/F} : B \to F$ for the reduced norm whereas often this denotes the full norm $(\Nm_{B/F})^2$ where $\Nm_{B/F}$, often written as $\mathrm{nrd}_{B/F}$, is the reduced norm).  Equivalently, since the definite quaternion algebra $\HH$ has only elements of positive norm, these are the elements $\gamma \in B$ so that $\iota_{\infty}(\gamma)$ is a sequence of real matrices with positive determinant.
\par
Via the embedding $\iota_{\infty} : B^+ \to \GL_2^+(\RR)^r$ we get an action $B^+ \acts \frH^r$ where $\frH \subset \CC$ is the upper half-plane. Quotients by special subgroups of $B^+$ will lead to Shimura varieties. The associated Shimura datum is for the algebraic group $G := \Res_{F/\Q} B^\times$, the functor $A \mapsto (B \times_{\Q} A)^\times$ on $\Q$-algebras, with the conjugacy class $X$ generated by the map $h : \SS \to G_{\RR}$ defined by
\[ \CC^\times \to \GL_2(\RR)^r \times \HH^{g-r} \quad \quad x + i y \mapsto \left( \begin{pmatrix}
    x & y
    \\
    -y & x
\end{pmatrix}, \dots, \begin{pmatrix}
    x & y
    \\
    -y & x
\end{pmatrix}, 1, \dots, 1 \right).  \]
In general, the Shimura datum $(G, X)$ is of abelian type but not necessarily of Hodge type \cite{milne:canonical_models}. Fix a maximal order $\cQ \subset B$ and set $\Gamma(\cQ, 1) := \cQ^\times \cap B^+$. The action $\Gamma(\cQ, 1) \acts \frH^r$ is totally discontinuous and cocompact as long as $B$ is not split over $F$. For a finite index subgroup $\Gamma \subset \Gamma(\cQ, 1)$ the quotients $Y = \frH^r / \Gamma$ are the complex points of the associated Shimura variety. For example:
\begin{enumerate}
    \item when $B = M_2(F)$ then $\cQ = M_2(\cO_F)$ and $Y$ is a Hilbert-Blumenthal modular variety which parametrizes abelian $2g$-folds with real multiplication by $\cO_F$
    \item $F = \Q$ and $B$ is a non-split indefinite quaternion algebra with discriminant $d$ then $Y$ is a Shimura curve parameterizing abelian surfaces with QM by $B$
    \item likewise, when $B$ is non-split over $F$ and $r = g$ then $Y$ is an honest moduli space parameterizing abelian $2g$-folds with QM by $\cQ$
    \item when $r = 1$ these $Y$ are the \textit{Shimura curves} but for $g > 1$ these are not honest moduli spaces of abelian varieties.
\end{enumerate}

These moduli interpretations make sense over a number field, leading to \textit{canonical models} for the modular Shimura varieties over this field. As in Shimura's original article, we focus here on the case $r = 1$. Even when $g > 1$ and the Shimura curve is not an honest moduli space, Shimura was able to construct canonical models over number fields by using isogenies to other moduli spaces. For an ideal $\frJ \subset \cO_F$, Shimura defines \textit{congruence subgroups}
\[ \Gamma(\cQ, \frJ) := \{ \gamma \in \cQ^\times \cap B^+ \mid \gamma - 1 \in \frJ \cQ \} \]
and he produces a canonical model for $\frH / \Gamma(\cQ, \frJ)$ as follows.

\begin{theorem}\cite[(3.2) Main Theorem I and (3.3)]{shimura:curves} \label{thm:canonical_models}
Let $F^{(\frJ \infty)}$ be the narrow ray class field of $\frJ$. Fix an embedding $F^{(\frJ \infty)} \embed \CC$ compatible with $\nu_1 : F \embed \RR$. Then there exists a curve $C$ defined over $F^{(\frJ \infty)}$ and a holomorphic map $\varphi : \frH \to C_{\CC}$ such that
\begin{enumerate}
    \item $\varphi$ descends to a biholomorphism $\frH  /\Gamma(\cQ, \frJ) \iso C_{\CC}$
    \item For any $K/F$ totally imaginary quadratic extension so that $B \ot_F K \cong M_2(K)$ and embedding $K \embed \CC$ compatible with $\nu_1 : F \embed \RR$ and any $F$-linear embedding $f : K \embed B$ so that $f(\cO_K) \subset \cQ$, let $z$ be the fixed point of $f(K) \acts \frH$; then $K^{(\frJ' \infty)}$ equals the compositum $K \cdot F^{(\frJ \infty)}(\varphi(z)) = K^{(\frJ' \infty)}$ where $\frJ' := \frJ \cO_K$ and $F^{(\frJ \infty)}(\varphi(z))$ is the field of definition of the point $\varphi(z) \in C(\ol{\Q})$.
\end{enumerate}
Furthermore, let $(C, \varphi)$ and $(C', \varphi')$ be two such pairs. Then there exists an isomorphism $j : C \to C'$ defined over $F^{(\frJ \infty)}$ such that $\varphi' = j_{\CC} \circ \varphi$.
\end{theorem}

Due to the uniqueness property, we call such a pair $(C, \varphi)$ a \textit{canonical model} and denote $C$ by $\X(\cQ, \frJ)$. This uniqueness property has the following consequence for automorphisms of $\X(\cQ, \frJ)$. Note that $\Gamma(\cQ, \frJ) \triangleleft \Gamma(\cQ, 1)$ so there is a well-defined action $\Gamma(\cQ, 1) \acts \frH/\Gamma(\cQ, \frJ)$ which induces an action on $\X(\cQ, \frJ)_{\CC}$ via $\varphi$.

\begin{lemma} \label{lem:automorphisms_defined_over_reflex}
$\Gamma(\cQ, 1)$ acts on $\X(\cQ, \frJ)$ via automorphisms \textit{defined over} $F^{(\frJ \infty)}$.
\end{lemma}

\begin{proof}
For any $\alpha \in \cQ^\times \cap B^+$, we can conjugate $f : K \to B$ to get $f^{\alpha} : K \to B$ with image $\alpha f(K) \alpha^{-1}$ and $f^{\alpha}(\cO_K) = \alpha f(\cO_K) \alpha^{-1} \subset \cQ$ since $\alpha \in \cQ^\times$ and $f(\cO_K) \subset \cQ$. Thus, if $z \in \frH$ is the fixed point of $f(K)$ then $\alpha \cdot z$ is fixed by $f^{\alpha}(K)$, so the hypothesis says that $K \cdot F^{(\frJ \infty)}(\varphi(\alpha \cdot z)) = K^{(\frJ' \infty)}$. Let $\varphi' := \varphi \circ \alpha$ viewing $\alpha$ as its action on $\frH$. Then the pair $(C, \varphi')$ is a canonical model because $\varphi(z') = \varphi(\alpha \cdot z)$ also satisfies $K \cdot F^{(\frJ \infty)}(\varphi(z)) = K^{(\frJ' \infty)}$ for any choice of $K/F$. Therefore, there exists $j : C \to C$ defined over $F^{(\frJ \infty)}$ so that $\varphi' = j_{\CC} \circ \varphi$, meaning exactly that the automorphism of $\X(\cQ, \frJ)$ induced by $\alpha$ is defined over $F^{(\frJ \infty)}$.
\end{proof}

\subsection{Congruence subgroups}

At split places of $B$ we can also define other congruence subgroups in analogy with the modular curves $X_0(p)$ and $X_1(p)$. Then we can fix an isomorphism:
\[ \iota_{\p} : B \ot_F F_{\p} \iso M_2(F_{\p}) \]
with the property that it restricts to an isomorphism $\iota_{\p} : \cQ \ot_{\cO_F} \cO_{F_\p} \iso M_2(\cO_{\p})$ (see \cite[Lemma~10.4.3]{voight:quaternion_algebras}).
Hence, reduction mod $\p^{\ell}$ gives an isomorphism
\[ \iota_{\p} : \cQ / \p \cQ \iso M_2(\kappa_{\p}) \]
where $\kappa_{\p} = \cO_F / \p$ is the residue field. This defines a reduction map
\[ \iota_{\p} : \Gamma(\cQ, 1) \to \GL_2(\kappa_{\p}) \]
whose kernel is exactly $\Gamma(\cQ, \p)$. More generally, if $\frJ \subset \cO_F$ is any ideal whose prime factorization consists entirely of primes splitting $B$, then we can fix a reduction map
\[ \iota_{\frJ} : \Gamma(\cQ, 1) \to \GL_2(\cO_F / \frJ).  \]
Taking the preimage of the upper Borel (resp.\ its unipotent radical, the upper triangular matrices) gives congruence subgroups $\Gamma_0(\cQ, \frJ)$ (resp.\ $\Gamma_1(\cQ, \frJ)$). Beware: this is not consistent with Shimura's use of the notation $\Gamma_1$. The groups $\Gamma_0(\cQ, \frJ)$ are the totally positive units in an \textit{Eichler order} $\cQ_0(\frJ) \subset \cQ$ of level $\frJ$ cf.\ \cite[\S3]{Voight:explicit_methods}.
\par
However, there are two caveats with these definitions: first $\iota_{\frJ} : \Gamma(\cQ, 1) \to \GL_2(\cO_F / \frJ)$ may not be surjective because of the restriction that elements live in $B^+$. Second, $\iota_{\infty}$ does not give a \textit{faithful} action of $\Gamma(\cQ, 1)$ on $\frH$ because the action factors through $\GL_2^+(\RR) \to \PGL_2^+(\RR)$. In order to understand the relationship between quotients by the various congruence subgroups, we need to precisely understand these failings.
\par
For any $\Gamma \subset \Gamma(\cQ, 1)$ denote by $\mathrm{P} \Gamma$ the quotient
\[ \mathrm{P} \Gamma := \frac{\Gamma}{\cO_F^\times \cap \Gamma} = \frac{\cO_F^{\times} \cdot \Gamma}{\cO_F^{\times}} \]
where $\cO_F^\times \embed \Gamma(\cQ, 1)$ are the scalar matrices. Since $\cQ^{\Nm = 1} \subset \Gamma(\cQ, 1)$ generates $\SL_2(\cO_{\p}) \subset \GL_2(\cO_{\p})$ it is clear that $\im{\iota_{\frJ}}$ contains $\SL_2(\cO_F / \frJ)$. Let
\[ \Gamma^1(\cQ, \frJ) := \Gamma(\cQ, \frJ)^{\Nm = 1} = \{ \gamma \in \cQ^{\Nm = 1} \mid \gamma - 1 \in \frJ \cQ \} \]
be the subgroup consisting of elements of $\Gamma(\cQ, \frJ)$ with reduced norm $1$. This is the kernel of $\iota_{\frJ}$ restricted to $\Gamma^1(\cQ, 1) := \cQ^{\Nm = 1}$.
We get a diagram of exact sequences:
\begin{center}
    \begin{tikzcd}
        & 0 \arrow[d] & 0 \arrow[d] & 0 \arrow[d] & 0 \arrow[d]
        \\
        0 \arrow[r] & L \arrow[d, hook] \arrow[r] & \Gamma^1(\cQ, \frJ) \arrow[d, hook] \arrow[r] & \mathrm{P} \Gamma(\cQ, \frJ) \arrow[r, "\Nm"] \arrow[d, hook] & U_{\p}^+ / U_\p^2\arrow[d, hook] &
        \\
        0 \arrow[r] & \{ \pm 1 \} \arrow[r] \arrow[d] & \Gamma^1(\cQ, 1) \arrow[r] \arrow[d, "\iota_{\frJ}"] & \mathrm{P} \Gamma(\cQ, 1) \arrow[r, "\Nm"] \arrow[d] & (\cO_F^\times)^+ / (\cO_F^{\times})^2 \arrow[d]
        \\
        0 \arrow[r] & \mu_2 \arrow[r] & \SL_2(\cO_F / \frJ) \arrow[r] & \PGL_2(\cO_F / \frJ) \arrow[r, "\det"] & (\cO_F / \frJ)^\times / (\cO_F / \frJ)^{\times 2} \arrow[r] & 0
    \end{tikzcd}
\end{center}
where $U_{\p} \subset \cO^\times_F$ is the group of $u \in \cO_F^{\times}$ such that $u \equiv 1 \mod \p$ and $L$ is $\{ \pm 1 \}$ if $2 \in \p$ and trivial otherwise. Therefore, when $(\cO_F^\times)^+ = (\cO_F^\times)^2$, for example, when the narrow class number of $F$ equals its class number, then the diagram shows that
\[ \Gamma^1(\cQ, \frJ) / L = \mathrm{P} \Gamma(\cQ, \frJ) \text{ and } \Gamma^1(\cQ, 1) / \{ \pm 1 \} = \mathrm{P} \Gamma(\cQ, 1) \text{ and } \im{\iota_{\frJ}} = \PSL_2(\cO_F / \frJ). \]
Therefore, in this case, $\X(\cQ, \frJ) \to \X(\cQ, 1)$ is a $\PSL_2(\cO_F / \frJ)$-cover. It has intermediate covers
\[ \X(\cQ, \frJ) \to \X_1(\cQ, \frJ) \to \X_0(\cQ, \frJ) \to \X(\cQ, 1) \]
which are quotients of $\X(\cQ, \frJ)$ by the Borel and its unipotent radical respectively.

Let $H \subset \GL_2(\cO_{\p})$ be an open subgroup for a prime $\p$ at which $B$ is split. The \textit{level} of $H$ is the smallest $\ell$ such that $H$ is the preimage of a subgroup $\ol{H} \subset \GL_2(\cO_F / \p^{\ell})$. Then we can define the Shimura curve $\X_H(\cQ)$ with $H$-level structure as the quotient of $\X(\cQ, \p^{\ell})$ by the subgroup $\Gamma_H(\cQ) := \Gamma(\cQ, 1) \cap H = \iota_{\p^{\ell}}^{-1}(\ol{H})$. By Lemma~\ref{lem:automorphisms_defined_over_reflex}, $\X_H(\cQ)$ also has a canonical model over $F^{(\frJ \infty)}$.

\subsection{Certain Hurwitz curves as Shimura curves}

Much of the favorable arithmetic of the genus 14 Hurwitz curves is due to their realization as Shimura curves for a particular quaternion algebra $B$ over $K_+ = \Q(\eta)$ where $\eta = \zeta_7 + \zeta_7^{-1} = 2 \cos{\tfrac{2\pi}{7}}$. The exposition here closely follows \S3.19 of \cite{shimura:curves} and Elkies' article \cite[\S4.4]{elkies:klein_quartic}. Choose as $B$ the following quaternion algebra:
\[ B = \left( \frac{\eta, \eta}{F} \right) := \frac{K_+ \langle i, j \rangle}{(i^2 - \eta, j^2 - \eta, ij + ji)}. \]
This is the unique quaternion algebra over $K_+$ that is non-split exactly over the two real places of $K_+$ where $\eta$ embeds negatively and split over all non-archimedean places. Let $v_{\infty} : K_+ \embed \RR$ denote the obvious real embedding sending $\eta \mapsto 2 \cos{\tfrac{2\pi}{7}}$ which splits $B$. We can thus fix an embedding $\iota_{\infty} : B \embed M_2(\RR)$ compatible with $v_{\infty}$ arising from a choice of isomorphism $B \ot_{K_+, v_{\infty}} \RR \iso M_2(\RR)$. We apply Shimura's construction from the previous section to this quaternion algebra. Define the following maximal order:
\[ \cQ_{\mathrm{Hur}} = \cO_{K_+}[i,j,j'] \subset B \]
where $j' := \tfrac{1}{2}(1 + \eta i + \tau j)$ and $\tau := 1 + \eta + \eta^2$. Since $K_+$ has narrow class number $1$, the previous section shows $\Gamma(\cQ,1) \onto \PSL_2(\cO_{K_+} / \frJ) \subset \PGL_2(\cO_{K_+} / \frJ)$.
\par
Miraculously,
\[ \mathrm{P} \Gamma(\cQ_{\Hur}, 1) = (\cQ_{\Hur})^{\Nm = 1} / \{ \pm 1 \} \cong \TT(2,3,7) = \left< g_1, g_2, g_3 \mid g_1^2, g_2^3, g_3^7, g_1 g_2 g_3 \right> \]
recovers the Hurwitz triangle group. Therefore, any finite-index normal subgroup $\Gamma \subset \Gamma(\cQ_{\Hur}, 1)$ without elliptic elements produces a cover $\X_{\Gamma} \to \X(\cQ_{\Hur},1) \cong \P^1$ ramified at exactly three points with orders $2,3,7$ and hence a Hurwitz curve with automorphism group $\rmP \Gamma(\cQ_{\Hur},1) / \rmP \Gamma$. % REVIEW: This sentence has a grammatical issue -- "any Hurwitz curve is uniformized ... is conjugate" lacks a main verb connecting the two clauses. Consider restructuring.
Conversely, any Hurwitz curve uniformized by $\frH$ with deck transformation group $\Gamma \subset \SL_2(\RR)$ (i.e., $\Gamma$ is a Fuchsian subgroup so that $\frH / \Gamma$ has $84(g-1)$ automorphisms) is conjugate in $\SL_2(\RR)$ to a normal subgroup of $\Gamma(\cQ_{\Hur}, 1)$ modulo the center.
\par
% REVIEW: "$\cQ_{\Hur}/\p \cong \SL_2(\kappa_\p)$" -- should this be $\cQ_{\Hur}/\p\cQ_{\Hur} \cong M_2(\kappa_\p)$ (as a matrix algebra), with the group statement being $\Gamma^1(\cQ_{\Hur},1) \onto \SL_2(\kappa_\p)$?
Because $B$ is non-split only at $\infty$ we know that $\cQ_{\Hur} / \p\cQ_{\Hur} \cong M_2(\kappa_{\p})$ for all primes $\p \subset \cO_{K_+}$. Therefore, % REVIEW: Should $\X(\cO, \p)$ be $\X(\cQ_{\Hur}, \p)$ to match the notation?
the Shimura curves $\X(\cQ_{\Hur}, \p)$ are Hurwitz curves with automorphism groups $\PSL_2(\kappa_{\p})$. By Theorem~\ref{thm:canonical_models}, this curve is defined over the ray class field $K_+^{(\p \infty)}$. For any prime $\p_{13}$ lying over $13$ this field $(K_+)^{(\p_{13} \infty)}$ is isomorphic to the field $K$ described in the main text (\cite[\href{https://www.lmfdb.org/NumberField/6.2.31213.1}{Number field 6.2.31213.1}]{lmfdb}). In fact, when $p$ is inert ($p \equiv \pm 2, \pm 3 \mod 7$) or ramified ($p = 7$) in $K_+$, the curves $\X^{\Hur}(\p)$ for the unique prime $\p$ over $p$ are defined over $\Q$ and when $p$ splits ($p \equiv \pm 1 \mod 7$) in $K_+$ the three curves $\X^{\Hur}(\p)$ for $\p \mid p$ are defined over $K_+$ and arise from the three embeddings $K_+ \embed \CC$ by \cite[Theorem~B]{voight:triangle_groups}
(cf. \cite[Proposition~5.1.2]{voight:thesis}). From now on, we make the abbreviation $\X^{\Hur}(\frJ) := \X(\cQ_{\Hur}, \frJ)$.

\begin{figure}
\begin{center}
    \begin{tabular}{c|c|c|c|c}
        $N(\p)$ & $\p$ & $m$ & $g$ & $G$
         \\
         \hline
         7 & $\p_7$ & 1 & 3 & $\PSL_2(\FF_7)$
         \\
         8 & $(2)$ & 1 & 7 & $\SL_2(\FF_8)$
         \\
         13 & $\p_{13}$ & 3 & 14 & $\PSL_2(\FF_{13})$
         \\
         27 & $(3)$ & 1 & 118 & $\PSL_2(\FF_{3^3})$
         \\
         29 & $\p_{29}$ & 3 & 146 & $\PSL_2(\FF_{29})$
    \end{tabular}
\end{center}
\caption{Prime ideals of $\cO_F$ ordered by norm. Here $m$ is the number of Galois conjugate prime ideals which all give distinct Hurwitz curves (see \cite[3.19.2]{shimura:curves}) each having genus $g$ and automorphism group $G$.}
\end{figure}

\begin{theorem}
All three Hurwitz curves of genus $14$ occur as $\X^{\Hur}(\p_{13})$ for the three primes $\p_{13}$ lying over $p = 13$. In particular, these three curves are Galois conjugate.
\end{theorem}

\begin{proof}
As tabulated in \cite[\href{https://www.lmfdb.org/HigherGenus/C/Aut/14.1092-25.0.2-3-7}{Higher Genus Family: 14.1092-25.0.2-3-7}]{lmfdb}, there are exactly three isomorphism classes of curves over $\CC$ of genus $14$ with automorphism group $\PSL_2(\FF_{13})$ and no other Hurwitz groups with order $1092$ exist. Furthermore, by \cite[3.19.2]{shimura:curves}, no two $\X^{\Hur}(\p_{13})$ are isomorphic over $\CC$ (base changed along the fixed embedding $K_+ \embed \RR$) for different choices of the prime $\p_{13}$ over $p$. Therefore, all three isomorphism classes of Hurwitz curves over $\CC$ appear.
\end{proof}

\begin{rmk}
Transitivity of the Galois action also follows from \cite[Theorem B (b)]{voight:triangle_groups} where it is shown that the $K_+$ is the field of moduli for $\PSL_2(\FF_{13})$-Galois $(2,3,7)$ \Belyi maps forgetting the isomorphism $\PSL_2(\FF_{13}) \iso \Aut(C)$. Note: given a Hurwitz curve, its $(2,3,7)$ \Belyi map is unique and canonically determined as the group scheme quotient $C \to C / \Aut(C)$.
\end{rmk}


\begin{theorem} \label{app:thm:isogeny_decomp_defined_over_K}
Let $\X^{\Hur}(\p_{13})$ be a genus $14$ Hurwitz curve. Then there is an elliptic curve $E$ defined over $K = (K_+)^{(\p_{13}\infty)}$ such that $\Jac(\X(\p_{13}))$ is isogenous to $E^{14}$ over $K$.
\end{theorem}

\begin{proof}
Let $C = \X^{\Hur}(\p_{13})$. Computing fixed points and comparing to genera of curves in the quotient lattice, we determine $H^1(C(\CC), \Q)$ as a $G$-representation to be $V_{14} \oplus V_{14}$ where $V_{14}$ is the $14$-dimensional irreducible $\Q$-representation of $G$ corresponding to the character \cite[\href{https://beta.lmfdb.org/Groups/Abstract/Qchar_table/1092.25?char_highlight=1092.25.14b}{1092.25.14b}]{lmfdb}. Moreover, the action of $\Q[G]$ on $V_{14}$ gives a surjection $\Q[G] \onto M_{14}(\Q)$ through which $\Q[G] \acts H^1(C(\CC), \Q)$ factors. Since $G \acts C$ is defined over $K$, the group algebra $\Q[G]$ acts by $K$-isogenies on $\Jac(C)$ so the idempotents in $M_{14}(\Q)$ give an isogeny decomposition
\[ \Jac(C) \sim E^{14} \]
defined over $K$.
\end{proof}

% REVIEW: Should "primes of $K_+$" be "primes of $K$"? The curve and the elliptic curve $E$ are defined over $K = (K_+)^{(\p_{13}\infty)}$, not $K_+$.
\begin{cor} \label{app:cor:supersingular_reduction}
$\X^{\Hur}(\p_{13})$ has supersingular reduction at infinitely many primes of $K$.
\end{cor}

\begin{proof}
Since $K$ has a real place, $E$ has infinitely many primes of supersingular reduction by \cite{elkies:supersingular}. Therefore, for infinitely many primes $\p \subset \cO_{K_+}$ there is a prime $\mathfrak{P} \subset \cO_{K}$ over $\p$ such that $E_{\mathfrak{P}}$ is supersingular. Therefore, whenever $C$ and $E$ have good reduction, we conclude that $C_{\mathfrak{P}}$ is supersingular since supersingularity is a  property of the Tate module of $\Jac(C)_{\mathfrak{P}}$ which is isomorphic to the Tate module of $E_{\mathfrak{P}}^{14}$. However, supersingularity is a geometric property, so $C_{\p}$ is supersingular as well.
\end{proof}

\section{Explicit computations in genera $3$ and $7$} \label{sec:genus3_7} \label{section:models}

(MODELS OF THE KLEIN QUARTIC)

{\color{red} TODO}

\section{Arithmetic of the first Hurwitz triplet}

The goal of this section is to explicitly determine the isogeny factors of the Jacobian of a genus $14$ Hurwitz curve. To do this, we need to use computations of quaternionic modular forms (or Hilbert modular forms via the Jacquet-Langlands correspondence) that are recorded in the LMFDB.

\subsection{Quaternionic modular forms}

In this section, we discuss modular forms invariant under action of congruence subgroups for a quaternion order $\cQ$. The relevance for us is that the Eichler--Shimura relation computes Frobenius traces on \etale cohomology in terms of Hecke eigenvalues.

\begin{defn}
For a matrix
\[ \gamma =
\begin{pmatrix}
a & b
\\
c & d
\end{pmatrix} \in \GL_2(\RR) \]
and $\tau \in \frH$ define the \textit{factor of automorphy}
\[ j(\gamma, \tau) := c \tau + d \in \CC \]
For a \textit{weight} vector $k = (k_1, \dots, k_n) \in (\Z^+)^r$, we define a right \textit{weight-k action} of $\Gamma(\cQ, 1)$ on functions $f : \frH^r \to \CC$ via
\[ (f|_k \gamma)(\tau) := f(\gamma \cdot \tau) \prod_{i = 1}^r (\det{\nu_i(\gamma)})^{k_i/2} j(\nu_i(\gamma), \tau_i)^{-k_i} \]
For $\Gamma \subset \Gamma(\cQ, 1)$ the space $M_k^B(\Gamma)$ of \textit{weight-k modular forms at level} $\Gamma$ is the $\CC$-vector space of holomorphic functions $f : \frH^r \to \CC$ invariant under the weight-$k$ action of $\Gamma$ and bounded at the cusps. The space $S_k^B(\Gamma)$ of \textit{weight-k cusp forms at level} $\Gamma$ is the subspace of functions vanishing at the cusps of $\frH^r / \Gamma$.
\end{defn}

In the case $B = M_2(F)$, $\cQ = M_2(\cO_F)$, and  $r = g$, we recover the spaces of Hilbert modular forms for the totally real field $F$. In that example, we drop the quaternion algebra ``$B$'' from the notation. We focus on the case $\Gamma = \Gamma_0(\cQ, \frJ)$ because these spaces have been computed in \cite{lmfdb} for small level. Write $S_k^B(\frJ) := S_k^B(\Gamma_0(\cQ, \frJ))$ and ditto for $M_k^B(\frJ)$. Parallel to the usual story, there is a Hecke algebra action on $S_k(\Gamma)$ now with Hecke operators $T_{\p}$ indexed by prime ideals $\p \subset \cO_F$. In the case when $B$ is non-split over $F$ it turns out there are no cusps (the corresponding subgroup is cocompact) so the condition at cusps is vacuous. In this case we say $S_k^B(\frJ) = M_k^B(\frJ)$ is the space of \textit{quaternionic cusp forms}.
\par
The case of parallel weight $2$, meaning $k = (2,\dots,2)$, is particularly important because such cusp forms define holomorphic symmetric $r$-forms on the Shimura variety $\frH^r / \Gamma$ (really the associated analytic quotient stack). In the case $r = 1$, this induces the \textit{Eichler--Shimura isomorphism}
\[ S_2^B(\frJ) \oplus \ol{S_2^B(\frJ)} \iso H^1(\X_0(\frJ), \CC). \]
Let $\TT_0(\frJ)$ be the quotient of the Hecke algebra acting faithfully on $S_2^B(\frJ)$. Then $H^1(\X_0(\frJ), \CC)$ becomes a free $\TT_0(\frJ)_{\CC}$-module of rank $2$. A Galois orbit of \textit{eigenforms} $f \in S_2^B(\frJ)$ for the Hecke operators is equivalent to an evaluation map $\ev_f : S_2^B(\frJ) \onto K_f$ sending $T_{\p} \mapsto a_{\p}(f)$ where $a_{\p}(f)$ is the eigenvalue of $T_{\p} f = a_{\p}(f) f$ and $K_f$ is the trace field generated by the $a_{\p}(f)$. The kernel defines a prime $\p_f = \ker{\ev_f}$ of the Hecke algebra. Since $\TT_0(\frJ)$ acts by correspondences, it acts on $\Jac(\X_0(\frJ))$ by isogenies so there is a quotient $A_f$ by $\p_f$ well-defined up to isogeny. Here $A_f$ is an abelian variety of dimension $[K_f : \Q]$ defined over $F$ \cite{zhang:heegner_points}. By above, its Tate module $V_{\ell}(A_f)$ is a $\Q_{\ell} \ot_{\Q} K_f$-module (whose action commutes with the Galois action) of rank $2$. The \textit{Eichler--Shimura relation} then says that for the Galois representation
\[ \rho_{A_f, \ell} : \Gal(\ol{F}/F) \to \GL_2(\Q_{\ell} \ot_{\Q} K_f) \]
and any prime $\p \subset \cO_F$ with $\p \nmid \ell \frJ$, the characteristic polynomial of $\rho_{A_f, \ell}(\Frob_{\p})$ is
\[ T^2 - a_{\p}(f) \, T + N(\p) \]
viewed as a polynomial over $\Q_{\ell} \ot_{\Q} K_f$ \cite[\S10.3]{carayol:representations}.
\par
The Jacquet--Langlands correspondence now relates the Hecke module structure for parallel weight modular forms in the quaternionic and Hilbert cases.

\begin{comment}
This construction of A_f is the ``Eichler--Shimura construction'' described for classical modular forms in ``Introduction to the arithmetic theory of automorphic functions'' or Brian's appendix to \href{https://wstein.org/papers/serre/ribet-stein.pdf}{Ribet-Stein}. However, in general for Hilbert-Modular forms the existence of $A_f$ with the right L-function is open (known as the ``Eichler--Shimura conjecture''). For Shimura curves, we have an actual Jacobian and hence can build the curve out of quotients by the Jacobian as in the classical case. However, in general, the way people prove cases of this conjecture appears to be using Jacquet-Langlands and then leverage the Jacobian of the associated Shimura curve. Of course, this only works when the associated quaternion algebra has exactly one real indefinite place.
\end{comment}

\begin{theorem}[Eichler--Shimizu--Jacquet--Langlands correspondence] \label{thm:jacquet_langlands}
Let $\frD \subset \cO_F$ be the discriminant of $B$ and $\frJ \subset \cO_F$ an ideal coprime to $\frD$. Then there is a Hecke module embedding
\[ S_2^B(\frJ) \embed S_2(\frD \frJ) \]
whose image consists of Hilbert cusp forms which are new at all $\p \mid \frD$.
\end{theorem}

See \cite[Proposition~2.12]{hida:abelian_varieties} (cf.\ \cite[Theorem~2.9]{Voight:computing_hecke} and \cite[Theorem~3.9]{Voight:explicit_methods}). In particular, when $\frD = (1)$, Theorem~\ref{thm:jacquet_langlands} yields an isomorphism $S_2^B(\frJ) \cong S_2(\frJ)$ as Hecke modules.

\subsection{Conjugation isomorphisms between Shimura curves}

Using conjugacy of certain Fuchsian groups, we can relate different Shimura curves at various levels. This will be important for explicit computations using the Jacquet-Langlands correspondence. The exposition closely follows the case of modular curves recorded in the appendix to \cite{zywina:images_of_galois} by John Voight.
\par
Let $H \subset \GL_2(\cO_{\p})$ be an open subgroup for a prime $\p$ at which $B$ is split. Assume $H$ is contained in the preimage of the upper Borel of $\GL_2(\kappa_{\p})$. Also assume that $\p$ has a totally positive generator $\pi \in \p$ (such as occurs when $F$ has narrow class number $1$). Consider the matrix
\[ \varpi := \begin{pmatrix}
0 & -1
\\
\pi & 0
\end{pmatrix} \in \GL_2^+(\cO_F) \]
This acts by conjugation in $\GL_2(F_{\p})$ as follows:
\[
\varpi
\begin{pmatrix}
a & b
\\
c & d
\end{pmatrix}
\varpi^{-1} =
\begin{pmatrix}
d & - \pi^{-1} c
\\
- \pi b & a
\end{pmatrix}
\]
Hence, if $c \in \p$ this produces a matrix in $\GL_2(\cO_{\p})$.
Therefore, there is a conjugation operation on the subgroups $H$ as follows:
\[ H^{\varpi} := \varpi H \varpi^{-1} = \left \{
\begin{pmatrix}
d & - c
\\
- \pi b & a
\end{pmatrix} \: \middle| \:
\begin{pmatrix}
a & b
\\
\pi c & d
\end{pmatrix}
\in H
\right \} \]
This conjugation action on subgroups induces nontrivial isomorphisms between Shimura curves at different levels. To prove this, we need an important result of Eichler.

\begin{lemma}[Eichler Approximation] \cite[Lemma~2.15]{shimura:curves}
Let $\cQ \subset B$ be a maximal order in a quaternion algebra over $F$ and $\mathfrak{I} \subset \cQ$ a two-sided ideal. Let $\beta \in \cO_F$ and $\gamma \in \cQ$ such that $\beta$ is positive at the non-split non-archimedean places of $B$ and $\Nm(\gamma) \equiv \beta \mod (\mathfrak{I} \cap \cO_F)$. Then there exists $\gamma' \in \cQ$ such that $\gamma' \equiv \gamma \mod \mathfrak{I}$ and $\Nm(\gamma') = \beta$.
\end{lemma}

\begin{prop} \label{prop:conjugation_isomorphism}
Let $H \subset \GL_2(\cO_{\p})$ be an open subgroup for a prime $\p$ at which $B$ is split contained in the preimage of the upper Borel of $\GL_2(\kappa_{\p})$. Then the Shimura curves $\X_H(\cQ)$ and $\X_{H^{\varpi}}(\cQ)$ are isomorphic as curves over $F^{(\frJ \infty)}$.
\end{prop}

\begin{proof}
We need to show that there is an automorphism of $\frH$ taking the action of $\Gamma_H(\cQ)$ to $\Gamma_{H^{\varpi}}(\cQ)$. A natural candidate is $\nu_1(\varpi)$ but it is not clear that this element actually does the job since conjugation by it may not preserve $B$. Instead, we will $\pi$-adically approximate $\varpi$ by an element of $\cQ$. Using the isomorphism $\iota_{\frJ} : \cQ / \p^{\ell} \cQ \iso M_2(\cO_F / \p^{\ell})$ we can find an element $\varpi_{\ell}' \in \cQ$ so that $\varpi_{\ell}' \equiv \varpi \mod \p^{\ell}$. Now we use Eichler's approximation theorem to find $\varpi_{\ell} \in \cQ$ such that $\Nm(\varpi_{\ell}) = \pi$ and $\varpi_{\ell} \equiv \varpi \mod \p^{\ell}$. It is clear that $(-)^{\varpi_{\ell}}$ preserves $B^+$. Since $\pi$ is totally positive, $\iota_{\infty}(\varpi_{\ell})$ acts on $\frH$. Now, if $H$ has level $\p^{\ell}$ then notice
\[ H^{\varpi} = H^{\varpi_{\ell}} \]
because $\varpi_{\ell} \gamma \varpi^{-1}_{\ell} \equiv \varpi \gamma \varpi^{-1} \mod \p^{\ell}$ and $H$ contains every element which is $1 \mod \p^{\ell}$. Now, $B^+$ is stable under $(-)^{\varpi_{\ell}}$ but $\cQ^{\times}$ is not, so we need to be careful in checking that $(-)^{\varpi_{\ell}}$ takes $\Gamma_H(\cQ)$ to $\Gamma_{H^{\varpi}}(\cQ)$. We do this at each local place using that
\[ \cQ = \bigcap_{\p' \subset \cO_F} (\cQ_{\p'} \cap B) \]
with the intersection taking place in $B$. Since $\Nm(\varpi_{\ell}) = \pi$ is a unit for $\p' \neq \p$, it preserves $\cQ_{\p'}$. At $\p$, an element of $\Gamma_H(\cQ)$ lands inside the maximal order $M_2(\cO_{\p})$. Therefore $(-)^{\varpi_{\ell}}$ sends $\Gamma_H(\cQ)$ inside $\cQ$. Its image must land inside $\cQ^{\times}$ since it is multiplicative and $\Gamma_H(\cQ)$ is a subgroup.
\end{proof}

For modular curves, Deligne--Rapoport prove that this isomorphism is defined over $\Q$ \cite[Propositions IV-3.16 and IV-3.19]{deligne:schemas_modules}. It is likely that the above isomorphism is defined over the reflex field $F^{(\p \infty)}$ but we will not need this.

\begin{example} \label{example:conjugation_p_squared_level}
For example, let $\p \subset \cO_F$ be a prime ideal generated by a totally positive element $\pi$ at which $B$ is split. The full congruence subgroup $\Gamma(\cQ, \p)$ is associated to the subgroup $H := \ker{(\GL_2(\cO_{\p}) \to \GL_2(\cO_F / \p)}$ which satisfies
\[ H^{\varpi} = \left \{
\begin{pmatrix}
d & - c
\\
- \pi b & a
\end{pmatrix} \: \middle| \:
\begin{pmatrix}
a & b
\\
\pi c & d
\end{pmatrix}
\in H
\right \} = \left \{ \begin{pmatrix}
1 + \pi d & - c
\\
- \pi^2 b & 1 + \pi a
\end{pmatrix} \: \middle| \:
a, b, c, d \in \cO_{\p}
\right\}  \]
which is the preimage of the intersection of the Borel in $\GL_2(\cO_F / \p^2)$ and the unipotent radical in $\GL_2(\cO_F / \p)$. Call this subgroup $H_{\mathrm{ub}}$ and $\X_{\mathrm{ub}}(\cQ, \p^2)$ its associated Shimura curve. Likewise, we can ask what larger subgroup $H'$ (hence also of level $\p$) corresponds to the full Borel $B(\p^2) \subset \GL_2(\cO_F / \p^2)$. This turns out to be the subgroup
\[ H_{\mathrm{sp}} = \left \{
\begin{pmatrix}
a & \pi b
\\
\pi c & d
\end{pmatrix}
\: \middle| \:
a, b, c, d \in \cO_{\p}
\right\}  \]
which is the preimage of the split Cartan subgroup of diagonal matrices in $\GL_2(\cO_F / \p)$. Indeed,
\[ H_{\mathrm{sp}}^{\varpi} = \left \{ \begin{pmatrix}
d & - c
\\
- \pi^2 b & a
\end{pmatrix} \: \middle| \:
a, b, c, d \in \cO_{\p}
\right\}  \]
is exactly the Borel in $\GL_2(\cO_F/\p^2)$. Therefore, Proposition~\ref{prop:conjugation_isomorphism} produces a diagram
\begin{center}
    \begin{tikzcd}
        \X_{\mathrm{ub}}(\cQ, \p^2) \arrow[d] \arrow[r, "\sim"] & \X(\cQ, \p) \arrow[d]
        \\
        \X_0(\cQ, \p^2) \arrow[r, "\sim"] & \X_{\mathrm{sp}}(\cQ, \p)
    \end{tikzcd}
\end{center}
\end{example}

\subsection{$\Jac(\X^{\Hur}(\p_{13}))$}

Here we compute the $j$-invariant of the elliptic curve $E$ appearing in the isogeny decomposition of $\Jac(\X^{\Hur}(\p_{13}))$. From Figure~\ref{fig:subgroups_genera} we know that $\X_{\mathrm{sp}}^{\Hur}(\p_{13})$, the quotient of $\X^{\Hur}(\p_{13})$ by the maximal torus $C_6$, is a genus $2$ curve. Since $\Jac(\X^{\Hur}(\p_{13})) \sim E^{14}$ over $K$ it suffices to compute $\Jac(\X_1^{\Hur}(\p_{13}))$. The conjugation isomorphism of Example~\ref{example:conjugation_p_squared_level} exhibits $\X_{\mathrm{sp}}^{\Hur}(\p_{13}) \iso \X_0^{\Hur}(\p_{13}^2)$. The Eichler--Shimizu--Jacquet--Langlands correspondence,
\[ S^B_2(\p_{13}^2) \cong S_2(\p_{13}^2) \]
expresses the Hecke action on $H^1(\X_0^{\Hur}(\p_{13}^2)_{\CC}, \QQ)$ in terms of Hilbert cusp forms of parallel weight $2$ for level $\Gamma_0(\p_{13}^2)$. Fixing the prime $\p_{13}$ over $13$, the space $S_2(\p_{13}^2)$ consists of a unique Galois orbit $f$ of Hilbert eigenforms \cite[\href{https://www.lmfdb.org/ModularForm/GL2/TotallyReal/3.3.49.1/holomorphic/3.3.49.1-169.4-a}{3.3.49.1-169.4-a}]{lmfdb}. The Hecke trace field is $K_f = \Q(\sqrt{3})$, hence $f$ corresponds to the abelian surface $A_f \cong \Jac(\X_0^{\Hur}(\p_{13}^2))$. We expect an elliptic curve $E$ in the isogeny decomposition over $K$ but the Hecke eigenvalues are not rational so the eigenform cannot correspond to an elliptic curve over $K_+$. Indeed, $K$ is the field of definition of $E$ and there is an isogeny $A_f \iso \Res^{K}_{K_+} E$. To find $E$ we analyze the Galois representation $\rho_{A_f, \ell}$ attached to the Tate module $V_{\ell}(A_f)$. Eichler-Shimura gives
\[ \tr{\rho_{A_f, \ell}(\Frob_{\p})} = a_{\p}(f) \]
for $\p \subset \cO_{K_+}$ not dividing $\ell \p_{13}$.
If $V_{\ell}(A_f)$ splits when restricted to $\Gal(\ol{K} / K)$, we can compute its traces of Frobenius directly. Consider $\frP \subset \cO_{K}$ lying over $\p \subset \cO_{K_+}$. If $\p$ is split or ramified then $\Frob_{\frP}$ and $\Frob_{\p}$ are conjugate in $\Gal(\ol{K_+} / K_+)$, so $\tr{\rho_{A_f, \ell}(\Frob_{\frP})} = a_{\p}(f)$. If $\p$ is inert, then $\Frob_{\frP}$ is conjugate to $\Frob_{\p}^2$ in $\Gal(\ol{K_+} / K_+)$ and therefore, by Cayley--Hamilton,
\[ \tr{\rho_{A_f, \ell}(\Frob_{\frP})} = (\tr{\rho_{A_f, \ell}(\Frob_{\p})})^2 - 2 \det{\rho_{A_f, \ell}(\Frob_{\p})} = a_{\p}(f)^2 - 2 N(\p). \]
As expected, every irrational value of $a_{\p}(f)$ appears for $\p$ inert in $K$ and $a_{\p}(f)^2 \in \ZZ$ meaning, in particular, the traces are the same for the two eigenforms exchanged by Galois. Hence, we get a unique list of rational traces $\tr{\rho_{A_f, \ell}(\Frob_{\frP})}$ for the primes $\frP \subset \cO_K$ corresponding to the decomposition $(A_f)_K \sim E^2$ also showing that $E$ is isogenous to its conjugate under $\Gal(K / K_+)$ (which also follows from the fact that $A_f$ has an isogeny action by $K_f = \Q(\sqrt{3})$). Since we know $E$ is defined over $K$ and has conductor dividing $\p_{13}^2$, since $\X_0^{\Hur}(\p_{13})$ has good reduction away from $\p_{13}$ by \cite{carayol:representations,kisin:integral_models}, there are finitely many potential elliptic curves. The built-in \textsc{Magma} function \texttt{EllipticCurveSearch} identifies $E$ and its conjugate from finitely many $\Frob_{\frP}$ traces (see the script \texttt{hilbert\_modular\_forms.m}). The minimal polynomial of $j(E) \in K$ is
\begin{align*}
    x^6 & - 14475 x^5 + 66476374 x^4 + 86360537888 x^3  - 1766214565832855 x^2
    \\
    & + 5634850735711464943 x - 5057260487992776789167.
\end{align*}
This elliptic curve has bad reduction exactly at the unique prime of $K$ lying over $\p_{13}$ which is also where $K / K_+$ is ramified. The first few supersingular primes of $E$ have norms $1091, 2339, 6551, 7643, \dots$.


\section{Genus $17$} \label{sec:genus17}

\subsection{Macbeath's construction}

In \cite{macbeath:hurwitz_theorem}, Macbeath shows that any Hurwitz curve $C$ of genus $g$ admits a canonical finite \etale cover $C_m \to C$ by another Hurwitz curve $C_m$ of genus
\[ g_m := m^{2g} (g-1) + 1. \]
The construction proceeds topologically as follows. Let $\pi_1(C) \to H_1(C, \Z)$ be the Hurewicz map. A characteristic subgroup of $\pi_1(C)$ will be a normal subgroup of $\pi_1(C) \subset \pi_1(\P^1 \sm \{ 0, 1, \infty \})$ and hence itself define a Hurwitz curve. One manifestly characteristic subgroup is $\Gamma_m \subset \pi_1(C)$ generated by all commutators and all $m$-th powers. This is the same as the kernel of $\pi_1(C) \to H_1(C, \Z / m \Z)$. Hence $\Gamma_m \subset \pi_1(C)$ defines a $(\Z / m \Z)^{2g}$-cover of $C$ which is also a Hurwitz curve. The Riemann-Hurwitz formula then gives
\[ g_m = g(C_m) = m^{2g}(g-1) + 1. \]

Applying this construction to the Klein quartic, for $m = 2$ we get a Hurwitz curve $C_2$ of genus $129$. However, $C_2 \to C$ has intermediate quotients. In general, there is an extension,
\[ 0 \to (\Z / m \Z)^{2g} \to \Aut(C_m) \to \Aut(C) \to 0 \]
and intermediate Galois quotients of $C_m \to C$ correspond to subgroups of $(\Z / m \Z)^{2g}$ stable under the $\Aut(C)$-action. % REVIEW: $(\Z/2\Z)^6$ has dimension 6 over $\FF_2$ but the Klein quartic has genus 3 so $H_1(C,\Z/2\Z) \cong (\Z/2\Z)^6$; confirm the exponent is correct here.
We will see that, for the Klein quartic, $(\Z / 2 \Z)^{6}$ splits as a direct sum of two irreducible $\Aut(X_3)$-representations over $\FF_2$. These give intermediate curves $(X_3)_2 \to X_{17,i} \to X_3$ for $i \in \{ 1,2 \}$ which are unramified $(\Z/2\Z)^3$-covers of $X_3$ and hence have genus $17$. These produce the only two Hurwitz curves of genus $17$.

\subsection{Arithmetic aspects of Macbeath's construction}

First we consider the general problem of producing geometrically irreducible \etale covers in arithmetic situations (or more generally over any non-algebraically closed field) from the geometric fundamental group. Let $X$ be a geometrically integral variety over $k$, fix a geometric point $\bar{x} : \Spec{\bar{k}} \to X$ and consider the homotopy exact sequence
\[ 1 \to \pi_1(X_{\bar{k}}) \to \pi_1(X) \to G_k \to 1 \]
Suppose we are given a finite index open subgroup $\Gamma \subset \pi_1(X_{\bar{k}})$ defining a finite \etale cover $\ol{Y} \to X_{\bar{k}}$. The question is when it can be extended to a map of exact sequences
\begin{center}
\begin{tikzcd}
1 \arrow[r] & \pi_1(X_{\bar{k}}) \arrow[r] & \pi_1(X) \arrow[r] & G_k \arrow[r] & 1
\\
1 \arrow[r] & \Gamma \arrow[u, hook] \arrow[r] & \wt{\Gamma} \arrow[u, hook] \arrow[r] & G_k \arrow[u, equals] \arrow[r] & 1
\end{tikzcd}
\end{center}
in which case $\wt{\Gamma} \subset \pi_1(X)$ defines a finite \etale cover $Y \to X$  pulling back to $\ol{Y} \to X_{\bar{k}}$. 

The existence of $\wt{\Gamma}$ is a purely group-theoretic question which we turn to next. We phrase the existence of $\wt{\Gamma}$ as a series of obstructions that must be overcome successively.
\begin{enumerate}
\item Let $N := N_{\pi_1(X)}(\Gamma)$ be the normalizer of $\Gamma$ in the arithmetic fundamental group $\pi_1(X)$. If $N \to G_k$ is not surjective then no $\wt{\Gamma}$ exists since $\wt{\Gamma}$ must live inside $N$. If $N \to G_k$ is not surjective, this froms the primary obstruction. 

\item Assuming $N \to G_k$ is surjective, now let $K := \ker{(N \onto G_k)}$. Next we must construct an extension
\begin{center}
\begin{tikzcd}
1 \arrow[r] & K \arrow[r] & N \arrow[r] & G_k \arrow[r] & 1
\\
1 \arrow[r] & \Gamma \arrow[u, hook] \arrow[r] & \wt{\Gamma} \arrow[u, hook] \arrow[r] & G_k \arrow[u, equals] \arrow[r] & 1
\end{tikzcd}
\end{center}
reducing the problem to the case that $\Gamma \subset \pi_1(X_{\bar{k}})$ is normal.  To do this, consider the quotient $Q := K / \Gamma$ and the induced diagram of exact sequences of groups,
\begin{center}
\begin{tikzcd}
1 \arrow[r] & Q \arrow[r] & E \arrow[r] & G_k \arrow[r] & 1
\\
1 \arrow[r] & K \arrow[r] & N \arrow[r] & G_k \arrow[r] \arrow[u, equals] & 1
\\
1 \arrow[r] & \Gamma \arrow[u, hook]
\end{tikzcd}
\end{center}
Here $E$ is constructed as the pushout of $K \onto Q$ and $K \embed N$ (PROVE THE SEQUENCE IS EXACT!). Then the existence of $\wt{\Gamma}$ is equivalent to the existence of $\Gamma_0 \subset E$ so that $\Gamma_0 \to G_k$ is an isomorphism. Indeed, then we take $\wt{\Gamma}$ to be the preimage of $\Gamma_0 \subset E$. This is the condition that the extension
\[ 1 \to Q \to E \to G_k \to 1 \]
is split on the right meaning the class $\ob \in H^2(G_k, Q)$ is trivial (NOTE ABOUT NONABELIAN COHOMOLGY?). 
\end{enumerate}

An important observation about the secondary obstruction $\ob \in H^2(G_k, Q)$ is as follows. Let $X_N \to X$ be the finite \etale cover corresponding to $N$, % REVIEW: These exact sequences are missing the rightmost term $G_k$; should read $1 \to K \to N \to G_k \to 1$ and $1 \to Q \to E \to G_k \to 1$.
if $X_N$ has a $k$-point then
\[ 1 \to K \to N \to G_k \to 1 \]
has a section inducing a section of
\[ 1 \to Q \to E \to G_k \to 1 \]
and hence the obstruction vanishes.

Let us specialize to the particular case when $\Gamma \subset \pi_1(X_{\bar{k}})$ is a characteristic subgroup. Then $N = \pi_1(X)$ and the first obstruction is trivial. Therefore, the obstruction to extending $\Gamma$ to $\wt{\Gamma}$ is a class 
\[ \ob \in H^2(G_k, Q). \]
For Macbeath's construction, $Q = H_{1,\et}(X_{\bar{k}}, \Z/m\Z) := \pi_1(X_{\bar{k}})^{\ab} \ot \Z / m \Z$ as a $G_k$-module and the obstruction lives in 
\[ H^2(G_k, H_{1,\et}(X_{\bar{k}}, \Z/m\Z)). \]
Furthermore, we know that if $X$ has a $k$-point then this obstruction vanishes. 

\subsection{The mod $2$ representation of the Klein Quartic} \label{sec:mod2_rep}

{\color{red} TODO}


{\color{red} from construction show that $A$ has full $2$-torsion over $\Q(\sqrt{-7})$. because its 2-homology has no Galois action by construction. THINK THIS IS FALSE}

In this section let $X$ be the unique $\Q$-model of the Klein quartic with full automorphism group defined over $\Q(\sqrt{-7})$ (see \S\ref{section:models}). Our goal is to compute
% REVIEW: Should this be $\GL_6(\FF_2)$ or $\GL(H_1(X_{\bar{\Q}}, \FF_2))$? The mod-2 representation of a genus-3 curve has dimension 6, not 2.
\[ \rho_{X,2} : G_{\Q} \to \GL_6(\FF_2) \]
We will see that $\rho_{X,2}|_{G_{\Q(\sqrt{-7})}} = \id$. This will then prove the following:

\begin{theorem}
The two Hurwitz curves of genus $17$ are defined over $\Q(\sqrt{-7})$ and are exchanged by the Galois action.
\end{theorem}

\begin{proof}
Notice that $X$ is isomorphic to the Klein quartic over $\Q(\zeta_7)$ and hence has a point over this field. Over {\color{red} FINISH}
\end{proof}

\subsection{The Jacobian of $X_{17}$}

% REVIEW: Should this be "over $\Q(\sqrt{-7})$" rather than "$\ol{\Q}(\sqrt{-7})$"? The previous theorem says these curves are defined over $\Q(\sqrt{-7})$.
In this section, we let $X_{17}$ denote either Hurwitz curve of genus $17$ over $\Q(\sqrt{-7})$.

The map $X_{17} \to X_3$ gives an isogeny decomposition $\Jac(X_{17}) \sim \Jac(X_3) \times \Prym(X_{17} / X_3)$ and % REVIEW: Should "$X_0(49)^3$" be "$\Jac(X_0(49))^3$" or "$E_0^3$" where $E_0$ is the elliptic curve $\Jac(X_0(49))$?
$\Jac(X_3)$ is a CM abelian $3$-fold isogenous to $E_0^3$. Therefore, it suffices to compute $\Prym(X_{17} / X_3)$. Luckily, we can further exploit the symmetries to reduce the computation to a Prym of a more tractable low genus cover.

From the description of $\Aut(X_3) \acts \FF_2^3$ we can easily find a presentation for $G_{1344} = \Aut(X_{17})$ a group of order $1344$. It is SmallGroup(1344, 11686) in the GAP library and an explicit presentation as a $(2,3,7)$-group can be extracted from this data \cite{Macbeath:eightfold, Conder:survey}.

We can compute the decomposition into irreducible $\Q$-reps of $G_{1344} \acts H^1(X_{17}, \Q)$ in a number of ways. One way is to consider all intermediate curves whose genera can be computed via the Riemann-Hurwitz formula and compare these numbers to the space of fixed points for the representation. We find that $H^1(X_{17}, \Q) \cong V_{14}^2 \times V_6$ where $V_{14}$ is an absolutely irreducible representation of dimension $14$ and $V_6$ is a representation of dimension $6$ factoring through $G_{1344} \to \Aut(X_3) \cong \PSL_2(\FF_7)$ which becomes a sum of two conjugate $3$-dimensional representations over $\Q(\sqrt{-7})$. This implies the $\Q$-algebra map $\QQ[G_{1344}] \to \End{H^1(X_{17}, \Q)}$ factors through an injection
\[ M_{14}(\Q) \times M_3(\Q(\sqrt{-7})) \embed \End{H^1(X_{17}, \Q)}. \] 
Each idempotent projects onto a $2$-dimensional subspace leading to an isogeny decomposition $\Jac(X_{17}) \sim E^{14} \times E_0^3$ where $E_0 = X_0(49)$ is an elliptic curve with CM by $\Q(\sqrt{-7})$. Then we can identify the $E_0^3$ factor with $\Jac(X_3)$ and the $E^{14}$ factor with $\Prym(X_{17}/X_3)$. Note this decomposition occurs only geometrically (at least over a field where the full group of automorphisms is defined).
\todo{here the point is we need to compute different $C_5$ them seperately}
To compute $E$ up to isogeny, we take an intermediate quotient $X_{17} \to C_5 \to X_3$ with $C_5$ a curve of genus $5$. Then $\Prym(C_5/X_3)$ is an abelian surface which is isogenous to $E^2$ since it maps to $\Prym(X_{17}/X_3)$. Since we have equations for $X_3$ we will be able to leverage explicit results on degree $2$ unramified coverings to compute equations for a curve $C_2$ such that $\Jac(C_2) \cong \Prym(C_5/X_3)$ and thus determine $E$.

First, we choose $C_5$ by quotienting by one of the $7 = 2^3 - 1$ hyperplanes of $\FF_2^3 \subset G_{1344}$ making up the kernel $\ker(\Aut(X_{17}) \to \Aut(X_3)) = \Aut(X_{17}/X_3)$. These are not $\Aut(X_3)$-invariant (since the representation is irreducible) so $C_5$ is not itself a Hurwitz curve. In fact, all $7$ choices of $C_5$ are permuted by $G_{1344}$. Furthermore, $\Prym(C_5/X_3) \to \Prym(X_{17}/X_3)$ span $\Prym(X_{17}/X_3)$ as we take different choices of $C_5$ {\color{red} HOW TO CHECK?}

\subsection{Prym varieties of degree $2$ covers of genus $3$ curves}

Given an \etale cover $\pi : D \to C$ of a genus $5$ curve over a genus $3$ curve, the \textit{Prym variety} 
\[ \Prym(D/C) :=  (\ker{\pi_* : \Jac(D) \to \Jac(C)})^{\circ} = \im{(\id - \iota_* : \Jac(D) \to \Jac(D))} \]
for $\iota$ the involution preserving $\pi$, is a principally polarized abelian variety of dimension $2$. Therefore, there is a uniquely determined -- by the Torelli theorem -- hyperelliptic curve $H$ such that $\Prym(D/C) \cong \Jac(H)$. If $\pi : D \to C$ is defined over a number field $K$ then $H$ and the isomorphism $\Prym(D/C) \cong \Jac(H)$ are all defined over $K$ (this holds because the map of stacks $\cM_2 \to \cA_2$ is a closed immesion \cite[Theorem 1.2]{aaron:torelli_map} -- a fact that is special to genus $2$). 
\par 
There is a well-known construction of $H$ explicitly in terms of equations for $\pi$ (see \cite{prym_genus_3}) which we review here. Suppose $\pi$ 

\todo{write down equations}

Now the question becomes, if we define a \etale cover of degree $2$ via a cohomology class $\alpha \in H^2_{\et}(\ol{C}, \FF_2)$ where $\ol{C}$ is the base change to the algebraic closure $\ol{K}$, how do we compute the equations for $H$. First, we should recall that even if $\alpha$ is $G_K$-invariant, we saw in the previous section there is an obstruction to $\alpha$ defining a cover defined over $K$. This obstruction comes from the 5-term sequence for the Leray/Hochschild-Serre spectral sequence,
\[ E^{p,q}_2 := H^p(G_K, H^q_{\et}(\ol{C}, \FF_2)) \implies H^{p+q}_{\et}(C, \FF_2) \]
computing absolute \etale cohomology. The 5-term sequence has the form
\[ 0 \to H^1(G_K, \FF_2) \to H^1_{\et}(C, \FF_2) \to H^1(\ol{C}, \FF_2)^{G_K} \to H^2(G_K, \FF_2) \to H^2_{\et}(C, \FF_2) \]
and therefore the obstruction to lifting along $H^1_{\et}(C, \FF_2) \to H^1(\ol{C}, \FF_2)^{G_K}$ lies in $H^2(G_K, \FF_2) = \Br(K)[2]$. We can rephrase this in terms of our previous obstructions. View $\alpha : \pi_1(\ol{C}) \to \FF_2$ and let $\Gamma := \ker{\alpha}$. If $\alpha$ is $G_K$-invariant then $\Gamma$ is stable under $G_K$-conjugation and $\Gamma \subset \pi_1(\ol{C})$ is normal so $N_{\pi_1(C)}(\Gamma) = \pi_1(C)$. Thus the primary obstruction vanishes and the secondary obstruction is a class $\ob \in H^2(G_K, Q)$ where $Q = \ker{(N \to G_K)} / \Gamma = \pi_1(\ol{C})/\Gamma = \FF_2$. 
\par 
Suppose we have a choose an extension $K_{\alpha} / K$ killing the obstruction $\ob(\alpha) \in \Br(K)[2]$. Then there exists an \etale double cover $\pi : D \to C$ defined over $K_{\alpha}$ and hence a quadratic decomposition $Q_1 Q_3 - Q_2^2 = \lambda F$ of the quartic polynomial $F$ cutting out $C$. It is computationally very expensive to find quadratic decompositions by exhaustively searching over quadrics with algebraic coefficients until we find a complete set of $2$-torsion points. However, we can exploit the structure of theta characteristics and their relation to bitangent lines of $C$ to reduce the search to finding all bitangents and certain quadrics witnessing so called syzygies of bitangents, a much more computationally feasible task. 


\subsection{Theta characteristics and bitangents}

\begin{defn}
A \textit{theta characteristic} $\theta$ on a smooth curve $C$ is a line bundle such that $\theta^{\ot 2} \cong \omega_C$.
\end{defn}

Theta characteristics have a natural notion of parity meaning they come in two types \textit{even} and \textit{odd} determined by the pairity of $h^0(C, \theta)$. We can give a nice interpretation of this pairty in terms of a quadratic form on $J(\ol{K})[2]$ where $J = \Jac(C)$. It is clear that the set of isomorphism classes of (geometric) theta characteristics is a $J(\ol{K})[2]$-torsor.
For each theta characteristic $\theta$ we can define a quadratic form on the $\FF_2$-vector space $J(\ol{K})[2]$.

\begin{lemma} \cite[Theorem~1.13]{harris:theta_chars}
The assignment
\[ q_{\theta} : \L \mapsto h^0(C, \L \ot \theta_0) - h^0(C, \theta) \mod 2 \]
defines a quadratic form on $J(\ol{K})[2]$ whose associated bilinear form is the Weil pairing composed with the unique isomorphism $\mu_2 \iso \Z / 2 \Z$.
\end{lemma}

In fact, the association $\theta \mapsto q_{\theta}$ defines an bijection $\Theta(\ol{C}) \iso Q(J(\ol{K})[2], \inner{-}{-})$ where $Q(V, \inner{-}{-})$ is the space of quadratic forms on $V$ whose associated bilinear form is $\inner{-}{-}$. Furthermore, $Q(V) := Q(V, \inner{-}{-})$ naturally decomposes into two subsets $Q(V)^+$ and $Q(V)^{-}$ of quadratic forms of Arf invariant $0$ and $1$ respectively. These are the two orbits of $Q(V)$ under the action of the symplectic group $\Sp(V)$. This decomposition gives the parity of a theta characteristic because
\[ \Arf(q_{\theta}) = h^0(C, \theta) \mod 2 \]
by \cite[Theorem~4]{johnson:spin_structures}. Atiyah proved that there are exactly $2^{g-1} (2^g + 1)$ even and $2^{g-1}(2^g - 1)$ odd theta characteristics and that the parity of $\theta$ is a deformation invariant (unlike $h^0(C, \theta) \in \Z_{\ge 0}$) \cite{atiyah:spin_structures}. 
\par 
In the particular case of $g = 3$, there are $36$ even and $28$ odd theta characteristics. The $28$ odd theta characteristics correspond exactly to the $28$ bitangents of the canonical embedding $C \embed \P^2$. Indeed, an odd theta characteristic is an effective divisor $D$ of degree $2$ such that $2D \sim K_C$ thus $2D \in |\ell|$ is in the linear series of lines. Since $C$ is irrational, $h^0(C, \cO_C(D)) \le \deg{D} - 1 = 1$ there is a unique representation, $D = P + Q$ for points $P, Q \in C$. Then $2 P + 2 Q = C \cdot \ell$ for some line $\ell \subset \P^2$ meaning exactly that the line through $P, Q$ (resp.\ the tangent line through $P = Q$) is a bitangent (resp.\ hyperflex).


\subsection{Symmetric genus $2$ curves}

\subsection{$4$-torsion obstruction}
\todo{not $X_{17}$ need a different curve}
A remarkable coincidence that allows us to conclude $X_{17}$ has infinitely many primes of supersingular reduction occurs in the congruence classes satisfied by rational primes above which $E$ is supersingular. Recall that $E$ is associated to a weight $2$ newform $f \in S_2^{\text{new}}(112, \chi_f)$ in the sense that the abelian surface $A = \Res^{\Q(\sqrt{7})}_{\Q} E$ over $\Q$ satisfies $\rho_{A,\ell} \cong \rho_{f,\ell}$. The nebentypus character of $f$ is the quadratic character $\chi_f(\bullet) = \legendre{28}{\bullet}$.
Furthermore, the coefficient field of $f$ is $K_f = \Q(\zeta_3)$ in which $2$ is inert. Therefore, since the $\Q(\zeta_3)$-multiplication implies that $E$ is isogenous to its Galois conjugate, we see that $E$ is supersingular at $\p \mid p$ iff $A$ is supersingular at $p$ iff $a_p(f) = 0$.

\begin{thm} \label{thm:numerical_coincidence}
Let $p$ be an odd prime split in $\Q(\sqrt{-7})$. Then $a_p(f) \neq 0$.
\end{thm}

This congruence restriction on $p$ such that $a_p(f) = 0$ gives our main result.

\begin{cor}
$X_{17}$ has infinitely many primes of supersingular reduction.
\end{cor}

\begin{proof}
Since $E / \Q(\sqrt{7})$ has real $j$-invariant, Elkies' theorem implies it has infinitely many primes of supersingular reduction i.e.\ there are infinitely many rational primes for which $a_p(f) = 0$. By the theorem, each such $p$ must be inert in $\Q(\sqrt{-7})$ and thus is automatically a prime of supersingular reduction of the elliptic curve $E_0 / \Q$ since $E_0$ has CM by $\Q(\sqrt{-7})$. Since $\Jac(X_{17}) \sim E^{14} \times E_0^3$ geometrically this implies that every prime lying over the infinitely many rational primes with $a_p(f) = 0$ gives supersingular reduction for $X_{17}$.
\end{proof}

To prove the theorem, we consider the $4$-torsion Galois representation of $A$,
\begin{center}
    \begin{tikzcd}
    G_{\Q} \arrow[rr, bend left, "\bar{\rho}_{A,4}"] \arrow[r, "{\rho_{A,2}}"'] & \GSp_4(\Z_2) \arrow[r] & \GSp_4(\Z/4\Z) \arrow[r] & \GSp_4(\FF_2)
    \end{tikzcd}
\end{center}

\begin{proof}[Proof of Theorem~\ref{thm:numerical_coincidence}]
Recall that $\rho_{f, \ell}(\Frob_p)$ has characteristic polynomial
\[ X^2 - a_p(f) X + p \, \chi_f(p) \]
Therefore, if $a_p(f) = 0$ then $\rho_{f,2}(\Frob_p^2)$ is the scalar matrix $-p \, \chi_f(p)$. In particular, it is also scalar under the isomorphism $\rho_{f,2} \cong \rho_{A,2}$ and hence on $4$-torsion we see that $\bar{\rho}_{A,4}(\Frob_p^2)$ is also the scalar matrix  $-p \, \chi_f(p) \mod 4$.

We compute the number field $K_4$ trivializing the $4$-torsion $A[4]$ (as the splitting field of the $4$-division polynomial of $A$, computed in Magma). This is number field \cite[\href{https://www.lmfdb.org/NumberField/16.0.24759631762948096.2}{Number field 16.0.24759631762948096.2}]{lmfdb} in the LMFDB. There is a factorization,
\begin{center}
    \begin{tikzcd}
    G_{\Q} \arrow[d, two heads] \arrow[r, "\bar{\rho}_{A,4}"] & \GSp_4(\Z/4\Z)
    \\
    \Gal(K_4/\Q) \arrow[ru]
    \end{tikzcd}
\end{center}
$K_4$ is Galois over $\Q$ with Galois group $D_4 \times C_2$. There are two important quadratic subfields of $K_4$
\begin{center}
    \begin{tikzcd}
    & K_4 \arrow[rd, dash, "C_2^3"] \arrow[ld, dash, "D_4"'] \arrow[dd, dash, "D_4 \times C_2"']
    \\
    \Q(\sqrt{7}) \arrow[rd, dash] & & \Q(\sqrt{-7}) \arrow[ld, dash]
    \\
    & \Q
    \end{tikzcd}
\end{center}
where we have marked the extensions by their Galois groups. In fact, the only essential information is that $\Gal(K_4 / \Q(\sqrt{-7}))$ has exponent $2$ which follows from $A$ having full rational $2$-torsion over $\Q(\sqrt{-7})$ (\S\ref{sec:mod2_rep}). Indeed, from the diagram
\begin{center}
\begin{tikzcd}
& G_{\Q(\sqrt{-7})} \arrow[r, hook] \arrow[d, dashed] & G_{\Q} \arrow[d, "\bar{\rho}_{A,4}"] \arrow[rd, "\bar{\rho}_{A,2}"]
\\
0 \arrow[r] & (\Z / 2 \Z)^8 \arrow[r] & \GL_4(\Z/4\Z) \arrow[r] & \GL_4(\FF_2) \arrow[r] & 0
\end{tikzcd}
\end{center}
we see that $\Gal(K_4 / \Q(\sqrt{-7})) \embed (\Z / 2 \Z)^8$ proving that it has exponent $2$.

Now suppose that $p$ is split in $\Q(\sqrt{-7})$ this means $\Frob_p \in G_{\Q(\sqrt{-7})} \subset G_{\Q}$ identified with $\Frob_{\p}$ for either $\p \mid p$ in $\cO_{\Q(\sqrt{-7})}$ (these are conjugate in $G_{\Q}$ but not in $G_{\Q(\sqrt{-7})}$). However, $\Gal(K_4/\Q(\sqrt{-7}))$ has exponent $2$ and therefore $\Frob_p^2$ is trivial in $\Gal(K_4 / \Q)$. In particular $\bar{\rho}_{A,4}(\Frob_p^2) = I$ but we also showed that $\rho_{A,2}(\Frob_p^2)$ is the scalar matrix $-p \, \chi_f(p)$. Hence we conclude that
\[ -p \, \chi_f(p) \equiv 1 \mod 4 \]
However, $p$ is split in $\Q(\sqrt{-7})$ so $\legendre{-7}{p} = 1$ and thus
\[ \chi_f(p) = \legendre{28}{p} = \legendre{4}{p} \legendre{-1}{p} \legendre{-7}{p} = \legendre{-1}{p} \]
because $p$ is odd. Therefore,
\[ p \equiv -\legendre{-1}{p} \mod 4 \]
which is impossible.
\end{proof}

\begin{rmk}
Warning: the $\Z_{2}$-lattice arising from the Tate module of $A$ does not agree with any lattice corresponding to an $\cO_{K_f, 2}$-integral structure of $\rho_{f,2}$. Indeed, otherwise $\bar{\rho}_{A,2}$ would factor through $\GL_2(\FF_4) \subset \GL_4(\FF_2)$ since $2$ is inert in $K_f = \Q(\zeta_3)$ but $A[2](\Q) = (\Z/2\Z)$ is not an $\FF_4$-vector space. This is witnessing the fact that $\zeta_3$ does not act on $A$, only some order inside $\Z[\zeta_3]$ does.
\end{rmk}

\ifSubfilesClassLoaded{%
  \bibliographystyle{alpha}
  \bibliography{refs}
}{}

\end{document}
